\documentclass[a4 paper, 11pt]{article}

\usepackage[ngerman]{babel}
\usepackage[T1]{fontenc}
\usepackage[utf8]{inputenc}

\usepackage{microtype}
\usepackage[top=3cm, bottom=2cm, inner=2cm, outer=2cm]{geometry}

\usepackage{parskip}
\usepackage{amsmath} 
\usepackage{amssymb} 
\usepackage{amsfonts}
\usepackage[locale=DE, per-mode=fraction, quotient-mode=fraction, separate-uncertainty = true]{siunitx}
\usepackage{nicefrac}
\usepackage{wasysym} %Durchmesserzeichen

\usepackage{textcomp}

\usepackage{placeins}

\usepackage{booktabs}
\usepackage{graphics}
\usepackage[rel]{overpic}

\usepackage{caption}

\usepackage[colorlinks, allcolors=black, citecolor=black!80]{hyperref}

\usepackage[all, defaultlines=3]{nowidow}

\usepackage{enumitem}

%\usepackage{comment} % Dient dazu, eine frühere Version bei Bedarf ein- oder ausblenden zu lassen.
%\includecomment{versiona}
%\excludecomment{versiona} 

%\usepackage{color} %farbige Schrift

%\definecolor{new}{RGB}{50, 150, 50} %green
%\definecolor{new}{RGB}{0, 0, 0} %black
%\definecolor{box}{RGB}{50, 150, 50} %green
%\definecolor{old}{RGB}{255, 69, 0} %orange
%\definecolor{move}{RGB}{50,100,255} %blue

%\usepackage{showframe}
%\usepackage{lipsum}

\usepackage{fancyhdr}

\fancyhead[L]{\textsc{Leitfaden} der Maultaschen Tübingen}
%\fancyhead[C]{}
\fancyhead[R]{Version vom \today}

%\renewcommand{\d}{\mathrm{d}}
%\renewcommand{\thesubsubsection}{\alph{subsubsection})}

\title{{\normalsize Ultimate Frisbee in Tübingen} \\ \textbf{\textsc{Leitfaden} der Maultaschen Tübingen}}
\author{Erste Auflage}

\date{Version vom \today \\ {\footnotesize zusammengetippt von Sonja Lorenz\footnote{\href{mailto:so.lorenz@web.de}{so.lorenz@web.de}}}}


\begin{document}
	
\pagestyle{fancy} % Damit wird Kopfzeile gedruckt.

\maketitle

\vspace{1cm}

\thispagestyle{empty}

\subsubsection*{Sinn und Zweck dieser Sache}
Dieser Leitfaden regelt und informiert über die Strukturen und Abläufe innerhalb des Ultimate Frisbee-Teams \glqq Maultaschen Tübingen\grqq.
Er dient damit unter anderem der Planung und Durchführung von Teamsitzungen, gibt allen Interessierten einen Überblick zum allgemeinen Status Quo und ist Grundlage für Entwicklungen und Entscheidungen in der Zukunft.

Im Anhang zu diesem Leitfaden sind zu bestimmten Themen weiterführende Details und Kommentare festgehalten. 

Aktuelle Informationen zu Ansprechpartnern, Trainingszeiten etc. finden sich auf der Maultaschen-Homepage unter \url{https://ultimate-tuebingen.de/}.

\vspace{2cm}

\subsubsection*{\glqq Spirit of the Game\grqq}
\begin{quote}
Ultimate relies upon a Spirit of the Game that places the responsibility for fair play on every player. There are no referees; the players are solely responsible for following and enforcing the rules, even at World Championship. Competitive play is encouraged, but never at the expense of respect between players, adherence to the rules and basic joy of play.

Spirit of the Game is the mindful behavior practiced by players worldwide prior to, during and after a game. It encompasses attitudes and skills such as good knowledge and application of the rules, fair-mindedness, safe play and spatial awareness, clear and calm communication, and a positive and respectful attitude towards teammates, fans, and opponents, in a mutual effort to protect the basic joy of play.
\end{quote}
\begin{flushright}
	\scriptsize{siehe wfdf.org, entnommen am 17. November 2019}
\end{flushright}

\newpage

\tableofcontents

\newpage

\section{Hard Facts}

\subsection{Wer sind die Maultaschen}
Die Maultaschen Tübingen sind ein Ultimate Frisbee-Team der Stadt Tübingen. 
Es ist dabei doppelt aufgestellt:
Zum einen bildet es die Abteilung \glqq Ultimate~Frisbee\grqq\ des Sportvereins SV03~Tübingen, zum anderen ist es Organisator der Ultimate-Kurse im Hochschulsportprogramm der Universität Tübingen.

\subsection{Ziele und Aufgaben}
Übergeordnetes Ziel ist die gemeinschaftliche Organisation und Förderung des Sports Ultimate~Frisbee und des damit verbundenen besonderen Sportsgeists (\glqq Spirit of the Game\grqq) im Raum Tübingen.
Das bedeutet insbesondere
\begin{itemize}
\item die Organisation und Durchführung von Training auf verschiedenen Leistungsniveaus,
\item das Miteinander als Team zu pflegen,
\item gemeinsam auf Turniere zu fahren und die Maultaschen Tübingen sowohl auf Spaß-Turnieren als auch bei offiziellen Turnieren (auf deutscher und internationaler Ebene) zu vertreten,
\item den Sport Ultimate Frisbee bekannter zu machen, selbst Turniere auszurichten und Nicht-Spieler\footnote{Es sei an dieser Stelle sicherheitshalber darauf hingewiesen, dass jegliche sprachliche Verwendung eines bestimmten Geschlechtes niemals als tatsächliche Einschränkung zu verstehen ist, sondern selbstverständlich immer alle Geschlechter miteinbezieht.} zu motivieren, Ultimate auszuprobieren,
\item Spaß am Spiel zu haben und weiterzugeben.
\end{itemize}

\subsection{Mannschaften}
Die Maultaschen Tübingen sind in den drei Divisionen Mixed, Open und Damen des Ultimate Frisbee Sportes vertreten.

Für die Abteilung Ultimate Frisbee des Vereins starten dabei die Mixed-Mannschaft unter dem Namen \glqq Maultaschen Tübingen\grqq, die Open-Mannschaft unter dem Namen \glqq THW\grqq\ sowie die Damen-Mannschaft als \glqq MissTÜrious\grqq.

Bei offiziellen Turnieren, die über Hochschulen oder Universitäten ausgerichtet werden (wie den deutschen Hochschulmeisterschaften), vertritt zudem die Unisport-Mannschaft \glqq Eberhard\grqq\ die Eberhard Karls Universität Tübingen.

\subsection{Logo und Farben}
Die Teamfarben der Maultaschen Tübingen sind gelb und lila.
Außerdem dient die folgende künstlerische Darstellung einer \glqq Maul-Tasche\grqq\ als Logo.
\begin{center}
	\includegraphics[width=0.15\textwidth]{logo_mt.png}
\end{center}


\section{Innere Struktur und Vorstand}

\subsection{Beschreibung der Organisationsstruktur}
Der Vorstand der Maultaschen Tübingen besteht aus fünf Personen, deren Aufgaben klar voneinander abgegrenzt sind (Details zur Aufgabenaufteilung stehen im Anhang). Diese fünf Posten sind:
\begin{itemize}
	\item Abteilungsvorstand (übernimmt die offizielle Abteilungsleitung im Verein)
	\item Hochschulsportvorstand
	\item Finanzvorstand
	\item Trainervorstand
	\item Spielervorstand
\end{itemize}
Die fünf Personen wählen zudem eine von ihnen zum Vorstandsvorsitzenden, der (bei Bedarf) die Koordination der Vorstandsarbeit federführend in die Hand nimmt.

Darüber hinaus gibt es vom Vorstand für bestimmte Bereiche/Aufgaben eingesetzte verantwortliche Personen, dazu zählen insbesondere:
\begin{itemize}
	\item Turnierkoordinatoren (via Onlinetool)
	\item Homepage-Beauftragter
	\item Facebook-Beauftragter
	\item Lichtmarkenverwalter
	\item Markenbring- und Schließbeauftragter fürs freie Spiel
	\item Scheibenverkäufer
	\item Pickup-Trikot-Verwalter
	\item Leitfaden-Verantwortlicher
	\item Maultaschen-Boxen-Mami
\end{itemize}
Außerdem regelmäßig Orgateams für die eigenen Turniere sowie Arbeitsgruppen zu aktuellen Themen wie beispielsweise für ein neues Scheibendesign, neue Trikots/Lifestylewear oder ähnliches.

Zudem ist es jedem Vorstand erlaubt, auf absehbare Zeit oder für konkrete Aufgaben einen oder mehrere von ihm bestimmte Stellvertreter zu benennen, die ihn bei seiner Arbeit unterstützen und ggf. auch vertreten (z.B. auf den monatlichen Abteilungsleitersitzungen des Vereins). Wird davon Gebrauch gemacht, sind die benannten Personen und ihre Befugnisse dem Team auf angemessene Art zu kommunizieren.

\subsection{Allgemeine Aufgaben, Rechte und Pflichten des Vorstands}
Ein guter Vorstand ist breit aufgestellt und kommuniziert harmonisch, um Agenden optimal voranzutreiben. Er entwickelt und pflegt das Gesamtkonzept Ultimate in Tübingen, übernimmt und koordiniert Aufgaben aus der internen Organisation und vertritt darüber hinaus die Anliegen der Maultaschen Tübingen nach außen.

Eine gute Dokumentation, offene und direkte Kommunikation mit dem Team und Transparenz helfen dabei, Missverständnissen vorzubeugen und eine aktive, fruchtbare Diskussionskultur zu pflegen.

Kleine Entscheidungen dürfen im Vorstand getroffen werden, ohne das Team darüber abstimmen zu lassen. Das schließt eine Finanzberechtigung für Einzelausgaben bis zu 500 Euro und innerhalb einer Amtszeit (vgl. Abschnitt~\ref{vorstandswahl}) von insgesamt 1500 Euro mit ein, solange es der Kontostand erlaubt. Wann das Team in eine Entscheidung mit einbezogen werden sollte, liegt im Ermessen der Vorstände. Diese sind daher aufgefordert, ihr Gespür für die Bedürfnisse und Wünsche im Team zu schärfen, um diesbezüglich gute Entscheidungen treffen zu können. Das bedeutet insbesondere, das Team rechtzeitig über entsprechende (Anschaffungs-)Überlegungen zu informieren und Rückmeldungen zu berücksichtigen.

Die Details zu den Verantwortlichkeiten der einzelnen Vorstandsposten sind im Anhang zu diesem Leitfaden gelistet.

\subsection{Wahl des Vorstands} \label{vorstandswahl}
Jedes Vereinsmitglied, das über 18 Jahre alt und damit voll geschäftsfähig ist, kann sich zur Wahl aufstellen lassen.

Abteilungsvorstand, Hochschulsportvorstand, Finanzvorstand und Spielervorstand werden jährlich in der Herbstsitzung (vgl. Abschnitt~\ref{teamsitzungen}) von allen stimmberechtigten Anwesenden gewählt. 
Stimmberechtigten, die sich aktuell selbst nicht als \glqq aktive\grqq\ Spieler wahrnehmen, bleibt es explizit freigestellt, sich bei der Wahl des Spielervorstands zu enthalten und sich so der Mehrheit der aktiven Spieler anzuschließen. Es ist dabei bewusst jeder Person selbst überlassen, sich als \glqq aktiv\grqq\ oder \glqq passiv\grqq\ einzuschätzen.

Der Trainervorstand wird in der Herbstsitzung ausschließlich von allen anwesenden, aktuell legitimierten Trainern (vgl. Abschnitt~\ref{trainermeetings} ) gewählt und ist durch diesen Prozess indirekt ebenfalls vom gesamten Team legitimiert.

Die Wahlen selbst laufen nach dem Schema einer normalen Abstimmung ab (vgl. Abschnitt~\ref{stimmberechtigung}), also offen per Handzeichen und es wird eine einfache Mehrheit\footnote{Eine einfache Mehrheit liegt vor, wenn eine Option mehr als die Hälfte der abgegebenen Stimmen ohne Enthaltungen erhalten hat.}
benötigt, um gewählt zu werden.

Die Vorstände sind anschließend angehalten, aus ihrem Kreise möglichst zeitnah einen Vorstandsvorsitzenden zu wählen.

Ist ein Vorstand aufgrund äußerer Umstände gezwungen, seinen Posten innerhalb einer Amtszeit vorzeitig abzugeben, so wird auf der nächstmöglichen Teamsitzung für den Rest der Amtszeit ein Nachfolger gewählt. Bis zur Wahl des Nachfolgers ist die Person offiziell noch im Amt, kann/darf aber ihre Aufgaben auf von ihr bestimmte Stellvertreter übertragen. 

Auf Antrag von mindestens drei Vereinsmitgliedern kann zudem ein Misstrauensvotum gegen einzelne Vorstände durchgeführt werden. Hierfür gelten die Bedingungen und Abläufe einer weitreichenden Entscheidung (vgl. Abschnitt~\ref{stimmberechtigung}).


\section{Regelmäßige Ereignisse und Abläufe} \label{regelmaessig}

\subsection{Teamsitzungen} \label{teamsitzungen}
\subsubsection{Sitzungstypen}
Es finden jährlich vier Teamsitzungen statt. Zwei große Teamsitzungen im Herbst und im Frühjahr (jeweils zu Semesterbeginn) mit Themenschwerpunkten bei Organisation/Finanzen/Vorstandwahl und zwei kleine Teamsitzungen dazwischen mit Themenschwerpunkten bei der Trainingsplanung für die kommende Saison. In allen Teamsitzungen können darüber hinaus je nach Bedarf noch aktuelle Themen diskutiert werden. 

Alle Interessierten, auch Gäste, sind herzlich eingeladen, aktiv an den Sitzungen teilzunehmen, es sind jedoch nur Vereinsmitglieder stimmberechtigt (siehe Abschnitt~\ref{stimmberechtigung}).

\subsubsection{Vorbereitung}
Jede Teamsitzung muss wenigstens zwei Wochen im Voraus vom jeweilig initiierenden Vorstandsmitglied (siehe Anhang) über den Emailverteiler angekündigt werden (eine zusätzliche Ankündigung per Whatsapp-Gruppe ist empfehlenswert). Diese Email enthält mindestens
\begin{itemize}
	\item Datum, Uhrzeit und Ort,
	\item die vorläufige Agenda (inkl. zugehöriger Dokumente),
	\item ggf. den Hinweis, dass man sich an den Finanzvorstand wenden möge, wenn man schon vor der Sitzung einen Blick auf den Finanzbericht werfen möchte
	\item und die Aufforderung, weitere Themenwünsche zeitnah an den Initiator rückzumelden, damit sie in die Agenda aufgenommen und falls notwendig noch separat angekündigt werden können. Insbesondere sollen nach Möglichkeit spontane Abstimmungen vermieden werden.
\end{itemize}

Zudem beauftragt der Initiator rechtzeitig im Voraus einen Protokollanten und, falls er die Sitzung nicht selber leiten möchte, einen Sitzungsleiter.

\subsubsection{Allgemeiner Ablauf}
Zu Beginn begrüßt der Sitzungsleiter alle Anwesenden, stellt die geplante Agenda vor, fragt nach, ob es noch Änderungsanträge gibt und prüft die Beschlussfähigkeit.
Beschlussfähigkeit liegt vor, wenn mindestens \SI{20}{\percent} aller Vereinsmitglieder anwesend sind.

Im Anschluss wird die Agenda Schritt für Schritt abgearbeitet, bis alles besprochen wurde (oder der Wirt den Laden schließt). Dabei wird (vom Sitzungsleiter) auf genügend Pausen und (von allen) auf konstruktive, prägnante Wortbeiträge geachtet, um die Diskussionen effizient voranzubringen. Der Sitzungsleiter hat eine moderierende Funktion und erteilt immer jeweils einer Person das Wort. 

Übliche Themen der großen Teamsitzungen sind dabei:
\begin{itemize}
	\item kurzer Rückblick auf die letzte Saison und Ankündigung der Pläne für die kommende (Trainingszeiten, Trainingstage, \dots)
	\item Überblick über den derzeitigen finanziellen Stand, im Frühjahr mit Bericht zur Kassenprüfung (wenn bereits erfolgt, sonst bei nächstmöglicher Gelegenheit)
	\item ggf. Entlastung\footnote{Der Verein bestätigt dem Vorstand durch die Entlastung, dass er die ihm übertragenen Aufgaben im Sinn des Vereins ordnungsgemäß erfüllt und die ihm anvertrauten Mittel des Vereins ordnungsgemäß verwaltet hat. [\dots] [D]em Vorstand [wird] durch die Entlastung im Nachhinein bestätigt, dass alles, was er mit den Mitteln des Vereins gemacht hat, in dessen Sinn war und durch diesen (nicht mehr durch den Vorstand persönlich) verantwortet wird.\linebreak \emph{Quelle:}~\url{https://de.wikipedia.org/wiki/Entlastung_(Recht)\#Entlastung_beim_Verein}, entnommen am 14.04.2020} der Vorstände, insb. Finanzvorstand
	\item Planung/Besprechung größerer Anschaffungen (z.B. neue Pavillons, Scheiben, \dots)
	\item im Herbst: Wahl der neuen Vorstände 
	\item im Frühjahr: Orga zum MTC
	\item im Herbst: Orgateam für das LALD finden
	\item Überblick zur Turnierplanung (Vorbereitungsturniere?)
\end{itemize}
\vspace{0.3cm}
Übliche Themen der kleinen Teamsitzungen sind dabei:
\begin{itemize}	
	\item Ankündigung der Trainingsplanung für die kommende Saison (Konzept, Trainingszeiten, \dots), d.h. Vorstellung der Ergebnisse der Trainersitzung (vgl. Abschnitt~\ref{trainermeetings})
	\item im Winter: Orga zum LALD	
	\item im Winter: Orgateam für den MTC finden
	\item im Winter: Legitimierung eines Kassenprüfers für die Jahresabrechnung
\end{itemize}

\subsubsection{Ablauf von Abstimmungen} \label{stimmberechtigung}
Steht (nach einer Diskussion) eine Abstimmung an, so wird zunächst unter Anleitung des Sitzungsleiters der genaue Wortlaut aller zur Wahl stehenden Optionen festgelegt und vom Protokollanten dokumentiert. Es sollte daraus insbesondere klar hervorgehen, ab wann der vorgesehene Beschluss in Kraft tritt. 

Wann bzw. ob eine Diskussion bereits abstimmungsfähig ist, entscheidet der Sitzungsleiter (in Absprache mit den anwesenden Personen).

Das weitere Vorgehen hängt nun davon ab, ob es sich um eine weitreichende Entscheidung handelt oder nicht. Eine Entscheidung ist dann weitreichend, wenn von mindestens drei stimmberechtigten Mitgliedern ein entsprechender (mündlicher) Antrag an den Vorstand gestellt wird. Idealerweise ist das bereits vor Beginn der Sitzung in der Vorbereitungsphase passiert, damit sich alle darauf einstellen können.

\emph{Weitreichende Entscheidung}
\nopagebreak

Ist eine weitreichende Entscheidung zu treffen, so findet die Abstimmung nicht sofort statt, um auch Mitgliedern ein Stimmrecht zu geben, die aus persönlichen Gründen der aktuellen Sitzung nicht beiwohnen können. 

In diesem Fall erstellt der Sitzungsleiter (oder eine von ihm beauftragte Person) nach der Sitzung eine anonyme Umfrage mit den gemeinsam formulierten Abstimmungsoptionen.
Sobald das Protokoll der Sitzung genehmigt ist (vgl. Abschnitt~\ref{protokolle}) kann dann in einem festen Zeitraum, mindestens jedoch eine Woche lang, online eine Stimme abgegeben werden.

Es sind dabei alle Vereinsmitglieder stimmberechtigt, die am Ende des Abstimmungszeitraums mindestens 16 Jahre alt sind. Jede Person hat \emph{eine} Stimme.

Am Ende der Abstimmung wird die Umfrage vom Ersteller ausgewertet und das Ergebnis publik gemacht (z.B. per Email und durch einen Nachtrag im Protokoll). Eine Option gilt dabei als angenommen, wenn sie mindestens eine \nicefrac{2}{3}-Mehrheit\footnote{Eine \nicefrac{2}{3}-Mehrheit liegt vor, wenn eine Option mindestens zwei Drittel der abgegebenen Stimmen ohne Enthaltungen erhalten hat.} erreicht hat und mindestens \SI{20}{\percent} der Vereinsmitglieder an der Abstimmung teilgenommen haben. Möglicherweise muss dazu in einem zweiten Schritt (auf die gleiche Art) eine Stichwahl zwischen den beiden Optionen durchgeführt werden, die in der vorherigen Runde die meisten Stimmen bekommen haben.\footnote{\label{fo:dreiOpt} Liegen drei (oder mehr) Optionen gleichauf, wird zwischen diesen nochmals abgestimmt.} Wird wiederholt keine \nicefrac{2}{3}-Mehrheit erreicht, wird die Entscheidung vertagt und in der nächsten Teamsitzung nochmals neu diskutiert.

\emph{Normale Abstimmung}
\nopagebreak

Handelt es sich nicht um eine weitreichende Entscheidung, so sind alle anwesenden Vereinsmitglieder (sowie alle Anwärter, von denen bereits eine unterschriebene Beitrittserklärung vorliegt) über 16 Jahre 
stimmberechtigt und haben jeweils \emph{eine} Stimme. Die Abstimmung erfolgt per Handzeichen.

Es werden nun vom Protokollanten nacheinander alle Optionen und zuletzt die Enthaltungen abgefragt. Erreicht dabei eine der Optionen eine einfache Mehrheit,
gilt sie als angenommen. Ist das nicht der Fall, so wird zwischen den beiden Optionen mit den meisten Stimmen eine Stichwahl durchgeführt.\textsuperscript{\ref{fo:dreiOpt}} Liegt eine Stimmengleichheit vor, geht das Thema nochmals in die Diskussion oder wird vom Sitzungsleiter auf die nächste Teamsitzung vertagt.

\subsubsection{Protokollgenehmigung} \label{protokolle}
Spätestens eine Woche nach der Sitzung ist vom Protokollanten das Protokoll teamintern zu veröffentlichen und über den Emailverteiler an alle zu verschicken.

Sobald das geschehen ist, dürfen eine Woche lang Einspruch erhoben und Änderungen beantragt werden. Passiert das und der Antrag wird vom Sitzungsleiter (in Absprache mit dem Vorstand) direkt angenommen, wiederholt sich der Vorgang mit der aktualisierten Version. Nach Ablauf der letzten Frist ist das Protokoll dann automatisch genehmigt.
	
Wird der Änderungsantrag nicht sofort angenommen, so wird dieser Teil (falls das möglich ist, ansonsten das ganze Protokoll) von der automatischen Genehmigung ausgenommen und in der nächsten Sitzung nochmals diskutiert und abgestimmt.

\subsection{Trainingsbetrieb}
\subsubsection{Trainersitzungen und -legitimation} \label{trainermeetings}
Zur Planung des zukünftigen Trainingsbetriebs wird vom Trainervorstand idealerweise im Voraus zu den kleinen Teamsitzungen im Winter und im Sommer eine Trainersitzung einberufen. Die Einladung erfolgt mit mindestens zwei Wochen Vorlauf über den Emailverteiler und richtet sich an alle, die an der Trainingsmitgestaltung und -entwicklung interessiert sind, unabhängig davon, ob sie sich vorstellen können, selbst ein Training zu leiten.

Der übliche Ablauf ist:
\begin{itemize}
	\item Rückblick auf die Trainings der vergangenen Saison und aktuelle Situation
	\item Formulierung der Ziele für die kommende Saison
	\item Wahl der Trainer durch alle Anwesenden per einfacher Mehrheit (soweit bereits Kandidaten zur Verfügung stehen)
\end{itemize}

Die Legitimation der Trainer durch das gesamte Team findet im Anschluss online über eine entsprechende Umfrage statt (analog zur Durchführung einer weitreichenden Entscheidung). Erhält ein Kandidat eine einfache Mehrheit, ist er in seinem Trainerposten bestätigt.

Können nicht alle Trainerposten direkt besetzt werden, ist es Aufgabe des Trainervorstandes, baldmöglichst entsprechende Freiwillige zu finden, die dann direkt das Legitimationsverfahren durchlaufen.

\subsubsection{Trainingsteilnahme}
An den Vereinstrainings dürfen prinzipiell alle Vereinsmitglieder teilnehmen. Um das Erreichen eines geplanten Trainingsziels, sinnvolles Training auf dem gewünschten Leistungsniveau und/oder eine gute Vorbereitung auf eine DM zu ermöglichen, können durch die zuständigen Trainer bzw. das zuständige Trainerteam jedoch geschlossene Trainingsgruppen benannt werden.

Die Anmeldung zu den Hochschulsporttrainings erfolgt über das Anmeldeverfahren der Universität und steht prinzipiell allen Interessierten offen.

\subsubsection{Freies Spiel}
Zusätzlich zu den verschiedenen Trainings findet einmal in der Woche das freie Spiel statt, das explizit allen Interessierten offen steht, egal \glqq ob sie erst 20 Minuten oder schon 20 Jahre Ultimate Frisbee spielen\grqq\ oder es einfach mal ausprobieren wollen. Es wird darauf geachtet, alle Spieler gleichermaßen einzubeziehen, es gibt viel Raum sich auszuprobieren und der Spaß am Spiel steht im Vordergrund.

\subsection{Turnierfahrten und Kader}
\subsubsection{Turnieranmeldung}
Alle Turniere, die nicht vom Deutschen Frisbeesportverband (dfv) oder als Hochschulmeisterschaften organisiert werden, sind prinzipiell offen für alle und werden über ein Onlinetool organisiert.

Dazu entscheiden die Turnier-Koordinatoren, ob ein (z.B. via Wurfpost angekündigtes) geplantes Turnier von den Maultaschen gespielt wird, schicken eine entsprechende Ankündigungs-Email über den Verteiler und fügen das Turnier im Online-Planungstool hinzu. Dort gilt das System \glqq first come, first serve\grqq, d.h. die ersten, die sich eintragen, dürfen das Turnier spielen, alle weiteren Eingetragenen rücken ggf. nach. Im Allgemeinen gilt die Anmeldeliste als voll, wenn sich jeweils doppelt so viele Damen oder Herren eingetragen haben, als maximal auf der Line benötigt werden.\footnote{Beispiel: Turnier mit 7v7 mit 1-2 Damen auf der Line $\Rightarrow$ max. 4 Damen und max. 12 Herren dürfen mitfahren.}
Es ist in Absprache mit den bereits eingetragenen Spielern aber natürlich möglich, dass noch mehr Leute mitfahren, wenn alle einverstanden sind.

Sobald sich genügend Teilnehmer eingetragen haben (im Zweifel mit den bereits Eingetragenen abzustimmen), kümmern sich die Turnier-Koordinatoren um die Anmeldung und die Überweisung der Teamfee. Außerdem finden sie unter den Teilnehmern einen Turnierpaten, dem sie alle relevanten Infos zum Turnier weiterleiten.
Es ist für die Planung daher essentiell, dass die eigenen Eintragungen als verbindlich wahrgenommen werden. Insbesondere kurzfristige Absagen sind für alle Beteiligten ärgerlich und daher zu vermeiden.

Falls ein Turnier als Vorbereitungsturnier für eine DM genutzt werden soll, muss dies so früh wie möglich gekennzeichnet werden. In diesem Fall haben Mitglieder aus dem Kader (vgl. Abschnitt~\ref{kader}) Vorrecht, am Turnier teilzunehmen. Andere Spieler wenden sich an den jeweils zuständigen Trainer, falls sie dennoch gerne am Turnier teilnehmen wollen.

Für die Anmeldungen bei den offiziellen Turnieren des dfv ist der Spielervorstand in Absprache mit den jeweiligen Trainern zuständig (außerdem muss darauf geachtet werden, auch die einzelnen Spieler rechtzeitig beim dfv anzumelden).
	
Die Anmeldung und Orga zur deutschen Hochschulmeisterschaft (DHM) erfolgt durch den Hochschulsportvorstand.

\subsubsection{Kader} \label{kader}
Für die offiziellen Turniere des dfv wird für jede der drei Divisionen von den jeweiligen Trainern ein Kader bestimmt. Dieser Kader vertritt die Maultaschen Tübingen für eine Saison bei Turnieren auf nationaler und internationaler Ebene.

Für den DHM-Kader ist der Hochschulsportvorstand zuständig. Er kann entweder selbst als Captain am Turnier teilnehmen und eine entsprechende Kaderwahl treffen, oder die Aufgabe teilweise oder vollständig an eine oder mehrere von ihm bestimmte Person(en) übertragen. Der Kader darf dabei nur aus Mitgliedern (Studenten und Mitarbeitern) der Universität Tübingen bestehen. Wenn rechtzeitig geplant, können jedoch eventuell Spielgemeinschaften mit anderen Hochschulen gebildet werden.


\section{Spezielle Regelungen} \label{speziell}

\subsection{Regelungen zur Turnierpauschale} \label{turnierpauschale}
Wurde die Teamfee für ein (gewöhnliches) Turnier, wie in den allermeisten Fällen üblich, aus dem Teamkonto bezahlt, so beteiligt sich daran jeder Mitspieler pauschal mit 10 Euro (bzw. 5 Euro bei Eintagesturnieren) unabhängig von der tatsächlichen Höhe der Teamfee. Der Betrag wird nach dem Turnier mit den Fahrtkosten und sonstigen Auslagen verrechnet (siehe Anhang). 

Ausgenommen davon sind Anfänger, die ihr erstes Turnier spielen.

War ein Spieler nicht die gesamte Spielzeit des Turniers dabei, kann außerdem der Finanzvorstand entscheiden, ob und inwiefern der Betrag im Einzelfall verringert wird.

\subsection{Regelungen zur Fahrtkostenabrechnung} \label{fahrtkostenabrechnung}
\subsubsection{Allgemeine Berechnung}
Üblicherweise geschieht die Anreise zu Turnieren mit privaten PKWs, die von einzelnen Spielern freiwillig zur Verfügung gestellt werden. Zur Umverteilung der dabei entstehenden Kosten werden (falls nötig für An- und Abreise separat) über eine Kilometerpauschale von 30\,\nicefrac{ct}{km} pro Auto die Gesamtkosten der Anfahrt bestimmt und gleichmäßig auf alle Mitspieler verteilt. Es ist den Fahrern dabei natürlich freigestellt, weniger zu berechnen, wenn sie das Auto z.B. von ihren Eltern geliehen haben und daher selbst nur die Benzinkosten tragen müssen. Reisen einzelne Spieler (z.B. weil nicht genügend Autoplätze zur Verfügung stehen) parallel dazu mit öffentlichen Verkehrsmitteln an, so werden die Kosten dafür zu den Gesamtkosten dazugerechnet und mitverteilt. Selbiges gilt für die Kosten eines eventuell organisierten Mietautos.

Sollte ein Mitspieler individuell an- und/oder abreisen, z.B. nicht von/nach Tübingen oder \glqq selbstverschuldet\grqq\ zu einem anderen Zeitpunkt, so trägt er die Kosten dafür selbstständig und ist für diese Fahrt aus der allgemeinen Kostenaufteilung ausgenommen.

Sonderfälle kann und darf es jedoch immer geben, daher sind die in diesem Abschnitt beschriebenen Regelungen nicht zwangsläufig bindend, sondern sollen im Speziellen unter den Turnierfahrern einvernehmlich geklärt werden.

\subsubsection{Offizielle Turniere}
Bei offiziellen Turnieren des deutschen Frisbeeverbands (dfv) wird vonseiten des Verbandes eine Fahrtkostenumverteilung organisiert, die unterschiedlich lange (und damit teure) Anfahrtswege zum Austragungsort ausgleichen soll. Der dadurch fällige oder erhaltene Betrag wird in der Turnierabrechnung gleichmäßig auf alle Mitspieler verteilt.

Zudem können bei der Stadt Tübingen für offizielle Turniere Fahrtkostenzuschüsse beantragt werden. Um auszugleichen, dass Teilnehmer weit entfernter Turniere deutlich höhere Fördersummen erhalten als Teilnehmer an lokalen Austragungsorten, werden grundsätzlich nur \SI{80}{\percent} der Fördersumme direkt an die Spieler weitergegeben, der Rest kommt dem gemeinsamen Konto zugute.
Da die Zuschüsse jedoch nur mit großer Zeitverzögerung (Größenordnung Monate bis über 1 Jahr) von der Stadt ausbezahlt werden und das MT-Konto es nicht leisten kann, die (erwarteten) Beträge so lange auszulegen, werden die Beträge den Spielern erst dann gutgeschrieben (siehe Anhang), wenn sie tatsächlich überwiesen wurden.

\subsection{Kommunikationsketten}
Bei Fragen, Ideen, Anregungen, Lob, Kritik, Unklarheiten und Wünschen jedweder allgemeiner (maultaschenbezogener) Art wendet man sich direkt an den zuständigen Vorstand, im Zweifel an den Spielervorstand. Falls es sich um ein spezielles Thema, Training, Turnier o.ä. handelt, kann auch zunächst der betreffende Verantwortliche/Trainer angesprochen werden, der das Anliegen dann ggf. an den Vorstand weitergibt. Ist man Anfänger oder neu im Team und kennt die entsprechenden Personen nicht, ist es natürlich erlaubt, dabei die Hilfe einer anderen Maultasche, z.B. eines (Anfänger-)Trainers, in Anspruch zu nehmen. 

Informationen zu den aktuellen Vorstandsmitgliedern und den von ihnen eingesetzten verantwortlichen Personen sowie den aktuellen Trainern sind auf der Maultaschen-Homepage unter \url{https://ultimate-tuebingen.de/} zu finden.


\section{Änderungen an diesem Leitfaden}
Änderungen an diesem Leitfaden sollten gut überlegt sein und müssen per Abstimmung 
beschlossen werden. Sie gelten mit Ausnahme von Änderungen an den Teilen~\ref{regelmaessig}, ohne~\ref{stimmberechtigung}, und~\ref{speziell} als weitreichende Entscheidungen und benötigen damit eine \nicefrac{2}{3}-Mehrheit.
Es sind hierbei die Inhalte der SV03-Satzung als Grundlage und Orientierungshilfe in den Entscheidungssprozess mit einzubeziehen.

Zur Durchführung einer Änderung sollte zum Zeitpunkt der Abstimmung allen Stimmberechtigten eine konkrete und vollständige neue Formulierung des betreffenden (neuen) Abschnitts vorliegen. Es muss zudem klar sein, ab wann diese Änderung in Kraft treten soll.

Ist ein Änderungsvorschlag angenommen, so ist es Aufgabe des vom Spielervorstand eingesetzten Leitfadenverantwortlichen, baldmöglichst eine aktualisierte Version des Leitfadens zu erstellen und für alle Maultaschen zugänglich zu machen.

\end{document}
